\chapter{Closed Loop Systems}
\label{chap:8}
This chapter covers the forward and inverse dynamics of rigid-body systems containing kinematic loops. The presence of kinematic loops brings a new level of complexity to the dynamics problem: new formulations are required; new problems arise; the systems exhibit new behaviours; and the inverse dynamics problem, in particular, changes its nature in the presence of kinematic loops. There are two main strategies for formulating the equations of motion for a closed-loop system: They are: 1. start with a set of unconstrained bodies, and apply all of the joint constraints simultaneously; or 2. start with a spanning tree of the system, and apply the loop-closure constraints. A third possibility, which will not be considered here, is to use the assembly method described in Section 7.5. Method 1 results in large, sparse matrix equations, and is considered briefly in Section 8.13. Method 2 is the best choice for typical closed-loop systems, and is the main topic of this chapter. Some of the special properties of closed-loop systems are discussed in Section 8.10; the use of a loop-closure function is covered in Section 8.11; and inverse dynamics is covered in Section 8.12. 8.1 Equations of Motion The equations of motion for general rigid-body systems, including closed-loop systems, were introduced in Chapter 3. In this section, we derive the equation of motion for a closed-loop system from that of its spanning tree. The procedure can be summarized as follows. 1. Formulate the equation of motion for a spanning tree of the closed-loop system.

2. Add terms to this equation representing the forces exerted on the tree by the loop joints. 3. Formulate a kinematic equation describing the motion constraints imposed on the tree by the loop joints. 4. Combine these two equations. If the spanning tree has n degrees of freedom, and the loop joints impose nc constraints on the tree, then this procedure results in a system of n + nc equations in n + nc unknowns. Consider a rigid-body system containing NB bodies, NJ joints and NL = NJ −NB kinematic loops. We assume that a spanning tree has been chosen, and that the bodies and joints have been numbered accordingly. The number of tree-joint variables, n, and the number of loop-closure constraints, nc, are given by the formulae n = NB X i=1 ni and nc = NJ X k=NB+1 nc k , (8.1) where ni is the number of freedoms allowed by tree joint i, and nc k is the number of constraints imposed by loop joint k. Observe that we are using the index variables i and k to range over tree-joint numbers and loop-joint numbers, respectively. More generally, we use the following index-variable convention in connection with kinematic loops: i and j range over tree-joint and/or body numbers (1 to NB); l ranges over loop numbers (1 to NL); and k identifies the loop joint that closes loop number l. With our standard numbering scheme, k and l are related by k = l + NB, and k therefore ranges from NB + 1 to NJ. From here on, wherever k and l appear together, it should be assumed that k = l + NB. The equation of motion for the spanning tree can be written H¨q + C = τ , (8.2) where ¨q and τ are vectors of tree-joint acceleration and force variables, respectively. Now, the only difference between the closed-loop system and its spanning tree is that the former contains the loop joints and the latter does not. Therefore, the equation of motion for the closed-loop system can be obtained from Eq. 8.2 by adding terms that account for the forces exerted on the tree by the loop joints. Thus, the equation of motion for the closed-loop system can be written H¨q + C = τ + τ c + τ a , (8.3) where τ c and τ a account for the constraint forces and active forces, respectively, produced by the loop joints. Active forces arise from springs, dampers, actuators, and the like, acting at the loop joints. If there are no such forces acting at a particular joint, then that joint is said to be passive; and if all of

the loop joints are passive, then τ a = 0. In this case, the dynamics calculation software can be simplified by omitting the code required to calculate τ a. τ a is assumed to be known; but τ c is unknown, and it can be expressed as a function of nc unknown constraint force variables (as in Eq. 8.5). Thus, Eq. 8.3 is a system of n equations in n + nc unknowns, and therefore needs to be supplemented with another nc equations before it can be solved. These equations come from the kinematic constraints. The loop joints impose a set of kinematic constraints on the tree, which can be collected into a single matrix equation of the form K¨q = k , (8.4) where K is an nc × n matrix. We say that this equation expresses the constraints at the acceleration level, on the grounds that it has been differentiated a sufficient number of times for the acceleration variables to appear. Expressions for K and k are derived in Section 8.2. As τ c is the force that imposes these constraints on the tree, it follows that τ c can be expressed in the form τ c = KTλ , (8.5) where λ is a vector of nc unknown constraint force variables. This step is covered in Section 8.4. (See also Eq. 3.15.) Equations 8.3, 8.4 and 8.5 can now be assembled into a single matrix equation, being the complete equation of motion for the closed-loop system:  H KT K 0   ¨q −λ  = τ −C + τ a k  . (8.6) This is a system of n+nc equations in n+nc unknowns. The coefficient matrix is symmetric, but not positive definite. If this matrix has full rank, then the equation can be solved for both ¨q and λ. If it does not have full rank, then some elements of λ will be indeterminate, but the equation can still be solved for ¨q. The rank of the coefficient matrix is n + rank(K). 8.2 Loop Constraint Equations Let us now formulate the coefficients of the loop constraint equation, Eq. 8.4. We start with the velocity across loop joint k, which is denoted vJk. This velocity is given by the equation vJk = vs(k) −vp(k) , (8.7) where p(k) and s(k) are the predecessor and successor bodies, respectively, of joint k. In the general case, joint k will impose a constraint of the form T T k vJk = T T k σk ,

where Tk and σk denote the constraint-force subspace and bias velocity, respectively, of joint k. (See Eq. 3.38.) However, for the sake of simplicity, we will assume that σk = 0 for all loop joints, so that the constraint equation simplifies to T T k vJk = 0 . (8.8) Combining Eqs. 8.7 and 8.8 gives a velocity constraint equation, T T k (vs(k) −vp(k)) = 0 , (8.9) which can be differentiated to obtain an acceleration constraint equation: T T k (as(k) −ap(k)) + ˙T T k (vs(k) −vp(k)) = 0 . (8.10) The velocity of any body in a kinematic tree can be expressed in terms of ˙q by an equation of the form vi = Ji ˙q , (8.11) where vi is the velocity of body i, and Ji is the Jacobian of body i. (See Example 4.2.) Ji is given by the equation Ji = ϵi1S1 ϵi2S2 · · · ϵiNBSNB  , (8.12) where ϵij takes the value 1 at every joint j that supports body i (i.e., the joint lies on the path between the body and the base), and zero otherwise. Expressed as an equation, ϵij =  1 if j ∈κ(i) 0 otherwise. Likewise, the acceleration of any body in a kinematic tree can be expressed in terms of ¨q by the derivative of Eq. 8.11: ai = Ji¨q + ˙Ji ˙q = Ji¨q + avp i . (8.13) avp i is the ‘velocity-product’ acceleration of body i, which is the acceleration it would have if all the tree-joint acceleration variables were zero. It is available as one of the by-products of calculating C. (See Section 8.6.) Combining Eqs. 8.10 and 8.13 gives T T k (Js(k) −Jp(k)) ¨q + T T k (avp s(k) −avp p(k)) + ˙T T k (vs(k) −vp(k)) = 0 , which we can simplify to T T k (Js(k) −Jp(k)) ¨q = kl (8.14) by defining kl = −T T k (avp s(k) −avp p(k)) −˙T T k (vs(k) −vp(k)) . (8.15)

Equation 8.14 is the acceleration constraint imposed on the spanning tree by loop joint k, and there is one such equation for each kinematic loop. Collecting these equations together, we get

T T NB+1 (Js(NB+1) −Jp(NB+1)) ... T T NJ (Js(NJ) −Jp(NJ))

¨q =

k1 ... kNL

. (8.16) On comparing this equation with Eq. 8.4, it can be seen that K =

T T NB+1 (Js(NB+1) −Jp(NB+1)) ... T T NJ (Js(NJ) −Jp(NJ))

and k =

k1 ... kNL

. (8.17) The expression Js(k) −Jp(k) defines a loop Jacobian—a matrix that maps ˙q to the velocity across a loop-closing joint. The loop Jacobian for loop l, denoted JLl, can be defined as follows: JLl = Js(k) −Jp(k) = εl1S1 εl2S2 · · · εlNBSNB  (8.18) where εlj = ϵs(k)j −ϵp(k)j. These numbers identify the tree joints that take part in each loop. If we define the root of loop l to be the highest-numbered body that is a common ancestor of both p(k) and s(k), then the set of tree joints that take part in loop l is the union of two subsets: those that connect the root to p(k), and those that connect the root to s(k). The value of εlj is −1 for each joint in the first subset, +1 for each joint in the second, and zero for all other joints. Using Eq. 8.18, we can formulate an alternative expression for K, which is the expression used to calculate it: K =

K11 · · · K1NB ... ... KNL1 · · · KNLNB

(8.19) where Klj = εljT T k Sj . (8.20) 8.3 Constraint Stabilization Equation 8.16 is theoretically correct, but is not stable during numerical integration. The problem is this: in principle, Eq. 8.16 behaves like a differential equation of the form ¨e = 0 ,

where e denotes a loop-closure position error; but in practice it behaves like ¨e = noise , where noise is a small-magnitude signal representing numerical errors, such as the truncation errors in the numerical integration process. Thus, the use of Eq. 8.16 will ensure that the acceleration errors are kept small, but there is nothing to stop an unbounded accumulation of position and velocity errors. The problem can be solved by adding a stabilization term, kstab, to Eq. 8.4 (and therefore also to Eq. 8.6), so that it now reads K¨q = k + kstab . (8.21) The purpose of kstab is to modify the behaviour of the constraint equation under numerical integration, so that it now behaves like an equation of the form ¨e + 2α ˙e + β2e = noise , where α > 0 and β̸ = 0. This method is due to Baumgarte, and is sometimes called Baumgarte stabilization (Baumgarte, 1972). The stabilization constants are usually chosen according to α = β = 1/Tstab , where Tstab is the desired time constant for the decay of position and velocity errors. To achieve this effect, kstab is defined as follows: kstab = −2α

T T NB+1vJNB+1 ... T T NJvJNJ

−β2

δ1 ... δNL

(8.22) where δl = jcalc( jtype(k), s(k),kXp(k),k) . (8.23) δl is a measure of the degree to which the position of the spanning tree violates the constraint imposed by loop joint k. It is a quantity of the form δl = T T k dl, where dl is a spatial vector measuring the position error around loop l. This vector depends only on the joint type and the transform that locates the joint’s successor frame relative to its predecessor frame. The latter is given by s(k),kXp(k),k = s(k),kXs(k) s(k)Xp(k) p(k)Xp(k),k = XS(l) s(k)Xp(k) XP(l)−1, (8.24) where XP and XS are the arrays of loop-joint predecessor and successor frame transforms stored in the system model. (See §4.2.) One problem with Baumgarte stabilization is that there is no hard-and-fast rule for selecting Tstab. A reasonable strategy is to ask how many snapshots

per second of the moving system you would need in order not to miss anything important, and choose a value of Tstab that is similar to, or a little shorter than, the period between snapshots. For example, if the system is a large industrial robot, then Tstab = 1 is too large, Tstab = 0.1 is reasonable, and Tstab = 0.01 is a bit small, but still reasonable. If Tstab is chosen too small, then the differential equations become unnecessarily stiff. The exact value of Tstab is not critical— changing it by a factor of 2 makes only a small difference. There is a temptation to choose Tstab as small as possible, so that position and velocity errors will decay as quickly as possible, in the belief that this will maximize the accuracy of the simulation. This is a bad strategy. The purpose of constraint stabilization is to achieve stability, not accuracy. If the simulation is not accurate enough, then the best way to improve it is to use a better integration method and/or a shorter integration time step. Calculating δ Let sXp be the coordinate transform from the predecessor frame to the successor frame of a loop joint, as determined by the position variables of the spanning tree. In general, this transform will not comply exactly with the motion constraint imposed by the loop joint. Therefore, let us factorize this matrix as follows: sXp = Xerr XJ , (8.25) where XJ complies exactly with the joint constraint, and Xerr represents the constraint error, which is assumed to be small. This factorization must be solved symbolically, using knowledge of the type of motion allowed by the joint, in order to obtain expressions for the elements of Xerr in terms of sXp. These expressions will be different for each joint type. Given that Xerr represents a small displacement, there exists a motion vector, d ∈M6, that satisfies 1 −d× ≃Xerr . (8.26) This vector is (approximately) the displacement from where the successor frame ought to be to where it actually is. Having obtained d, δ is then given by δ = T Td . (8.27) Formulae for δ for various common joint types are shown in Table 8.1. Note that a zero-DoF joint can serve as a universal loop joint, since any closed-loop system can be modified, by adding massless bodies and zero-DoF joints, in such a way that every loop-closing joint is a zero-DoF joint. Example 8.1 The formula in Table 8.1 for the cylindrical joint can be obtained as follows. By combining Eqs. 8.25 and 8.26, and choosing expressions for XJ and d that are appropriate for a cylindrical joint, we get the following

Zero DoF Revolute Prismatic Cylindrical 1 2

X23 −X32 X31 −X13 X12 −X21 X53 −X62 X61 −X43 X42 −X51

X23 −X13 X53 −X43 X41X21 + X42X22

1 2(X23 −X32) 1 2(X31 −X13) 1 2(X12 −X21) X53 −X43

X23 −X13 X53 −X43

Spherical

X51X31 + X52X32 + X53X33 X61X11 + X62X12 + X63X13 X41X21 + X42X22 + X43X23

\begin{table}[ht]\centering
\caption{Formulae for δ = jcalc(type, X) approximate equation: sXp ≃}
\label{table:8.1}
\end{table}
1 0 −d2 0 0 0 0 1 d1 0 0 0 d2 −d1 1 0 0 0 0 0 −d5 1 0 −d2 0 0 d4 0 1 d1 d5 −d4 0 d2 −d1 1

c s 0 0 0 0 −s c 0 0 0 0 0 0 1 0 0 0 0 z 0 c s 0 −z 0 0 −s c 0 0 0 0 0 0 1

, where c = cos(θ), s = sin(θ), and θ and z are the angular and linear displacements of the cylindrical joint (i.e., the joint position variables). Observe that we have set d3 = d6 = 0 as they are the coordinates of d in the directions of motion freedom of the joint. If we perform the multiplication, and examine the result, then we find that d1 = (sXp)23, d2 = −(sXp)13, and so on. Finally, δ = [d1 d2 d4 d5]T. 8.4 Loop Joint Forces Let fk denote the force transmitted across loop joint k in the closed-loop system. The effect of joint k on the spanning tree is therefore to exert a force of fk on body s(k) and a force of −fk on body p(k). Now, if a spatial force of f is applied to body i in a kinematic tree, then it has the same effect on the tree as a joint-space force of τ, where τ = J T i f . Therefore, the effect of joint k on the spanning tree is equivalent to a joint-space force of τ = (J T s(k) −J T p(k))fk = JT Ll fk .

Let us decompose fk into the sum of an active force, f a k , and a constraint force, f c k , and consider their effects separately. If τ a represents the net effect on the spanning tree of all the active forces at the loop joints, then τ a is given by τ a = NL X l=1 J T Ll f a k . (8.28) Likewise, if τ c represents the net effect of all the constraint forces at the loop joints, then τ c is given by τ c = NL X l=1 J T Ll f c k . (8.29) However, the constraint force at joint k can also be expressed in the form f c k = Tk λk , (8.30) where λk is a vector of nc k unknown constraint force variables (cf. Eq. 3.34), so τ c can also be expressed as τ c = NL X l=1 J T Ll Tk λk . (8.31) On comparing this equation with the definition of K in Eq. 8.17, it can be seen that τ c = KTλ (8.32) where λ =  λT NB+1 · · · λT NJ T . (8.33) As mentioned earlier, the active forces are assumed to be known. The best way to incorporate τ a into the equation of motion is to modify the algorithm for calculating C so that it calculates C −τ a instead. A suitable algorithm is presented in Section 8.6. 8.5 Solving the Equations of Motion Having formulated every term in the equations of motion, the next step is to solve them for the tree-joint accelerations. To recap, the equation of motion for the spanning tree, subject to the loop-closure forces, is H¨q + C = τ + KTλ + τ a ; (8.34) the loop-closure constraint equation is K¨q = k + kstab ; (8.35)

and the two can be combined into a single composite equation,  H KT K 0   ¨q −λ  =  τ −C + τ a k + kstab  . (8.36) In these equations, H is an n × n symmetric, positive-definite matrix, while K is a general nc × n matrix that may be rank-deficient. The coefficient matrix in Eq. 8.36 is therefore symmetric, and it has dimensions of (n + nc) × (n + nc); but it is not positive definite, and it has a rank of n + r, where r = rank(K), so it will be singular whenever r < nc. If r = nc, then both ¨q and λ are uniquely determined by Eq. 8.36. However, if r < nc, then nc −r of the constraints imposed on the spanning tree by the loop joints are linearly dependent on the other r constraints. As a result, some elements of λ will be indeterminate, but ¨q is still uniquely determined by the equation. Another problem that can arise if r < nc is that Eq. 8.35 can be slightly inconsistent; that is, k + kstab can lie slightly outside range(K). This problem is caused by numerical errors in the calculation of K, k and kstab, and by inexact satisfaction of the loop-closure constraints for positions and velocities. When Eq. 8.35 is slightly inconsistent, it is acceptable to solve it in a least-squares sense. There are various ways to solve the above equations for ¨q. Let us now examine some of them. In particular, let us consider the following three methods: 1. solve Eq. 8.36 directly; 2. solve Eq. 8.36 for λ, and use the result to obtain ¨q; or 3. solve Eq. 8.35 to express ¨q in terms of a vector of independent acceleration variables, and then solve Eq. 8.34. Method 1: Solve Eq. 8.36 Directly If the coefficient matrix in Eq. 8.36 is nonsingular, then this equation can be solved directly by a general-purpose linear equation solver, or a solver designed for symmetric, indefinite matrices. This is the simplest solution method, but not the most efficient. It is therefore appropriate whenever human effort is more important than computational efficiency. The coefficient matrix will be nonsingular if r = nc, or if nc −r linearly-dependent rows are pruned out of Eq. 8.35 before it is incorporated into Eq. 8.36. The latter is practical if one happens to know in advance which rows to remove. This pruning tactic amounts to reducing nc until nc = r. Pruning can also be used in conjunction with methods 2 and 3. Method 2: Solve for λ First By subtracting KH−1 times the first row from the second row in Eq. 8.36, we get an equation for λ of the form A λ = b , (8.37)

where A = KH−1KT (8.38) and b = k + kstab −KH−1(τ −C + τ a) . (8.39) The coefficient matrix, A, is symmetric and positive-semidefinite. Its size is nc × nc, and its rank is r. If r = nc then A is invertible, and Eq. 8.37 has the unique solution λ = A−1b . (8.40) If r < nc then A is singular, and Eq. 8.37 has an infinity of solutions. In this case, the general form of the solution is λ = A+ b + (1 −A+A)z , (8.41) where A+ is the pseudoinverse of A, z is an arbitrary vector, and the operator 1−A+A projects vectors onto the null space of A. If Eq. 8.37 is consistent, then Eq. 8.41 describes a family of exact solutions; otherwise, it describes a family of least-squares solutions. Equation 8.37 is consistent if and only if Eq. 8.35 is consistent. Having obtained a value for λ, the final step is to substitute it into Eq. 8.34, and solve the resulting equation for ¨q. Note that z has no effect on ¨q, because KT(1 −A+A) = 0 . One is therefore free to choose any convenient value for z, such as z = 0. The following procedure will calculate ¨q in an efficient manner. It uses the LTL factorization and back-substitution algorithms described in Section 6.5, and therefore exploits any branch-induced sparsity that may be present in H. It can easily be modified to use the LTDL factorization instead. 1. Factorize H into H = LTL. 2. Calculate τ ′ = τ −C + τ a. 3. Using back-substitution, calculate Y = L−TKT and z = L−Tτ ′. 4. Calculate A = Y TY and b = k + kstab −Y Tz. 5. Solve Aλ = b for λ by any appropriate method. 6. Solve H¨q = τ ′ + KTλ for ¨q using the factors from step 1. Method 3: Solve Eq. 8.35 First Suppose we have been given a general linear equation, Ax = b, in which A is an m × n matrix of rank r, and b ∈range(A) (i.e., the equation is consistent). The task is to solve it for x. If r = n then there will be only a single solution, which is given by the formula x = A+b. If r < n then there will be an infinite number of solutions, which are given by the formula x = Cy + d. In this

equation, y is an (n −r)-dimensional vector of unknowns, d is any one vector satisfying Ad = b, and C is a full-rank matrix having the property AC = 0. Each solution has a unique corresponding value of y, and vice versa (i.e., there is a 1 : 1 correspondence between the two). Let us apply this formula to Eq. 8.35. This equation has an nc×n coefficient matrix of rank r. If r = n then the equation has a unique solution, which is ¨q = K+(k + kstab); but if r < n then there are infinitely many solutions, all given by the formula ¨q = G ¨y + g . (8.42) In this equation, G is an n × (n −r) full-rank matrix having the property KG = 0 , (8.43) ¨y is an (n −r)-dimensional vector of independent acceleration variables, and g is any one vector satisfying Kg = k + kstab . (8.44) Typically, ¨y is chosen to be a subset of the elements of ¨q. It can therefore be thought of as a vector of independent joint acceleration variables for the closedloop system. An algorithm for calculating G and g is presented in Section 8.8. We now have an equation that expresses ¨q as a function of the unknown vector ¨y. The next step is therefore to find an expression for ¨y. Premultiplying Eq. 8.34 by GT causes λ to be eliminated, resulting in GTH ¨q = GT(τ −C + τ a) ; and substituting Eq. 8.42 into this equation produces GTHG ¨y = GT(τ −C + τ a −Hg) (8.45) (cf. Eq. 3.20). This equation can be regarded as the equation of motion for the closed-loop system, expressed in terms of the independent joint acceleration variables in ¨y. The matrix GTHG is symmetric and positive-definite, so Eq. 8.45 can be solved immediately for ¨y, giving ¨y = (GTHG)−1GT(τ −C + τ a −Hg) . (8.46) 8.6 Algorithm for C −τ a The best strategy for dealing with τ a is to choose the spanning tree such that every loop joint is passive. In this case, f a k = 0 for every loop joint, and so τ a = 0. The next best strategy is to modify the code for calculating C so that it calculates C −τ a instead.

Original Equations: vi = vλ(i) + Si ˙qi (v0 = 0) ai = aλ(i) + Si ¨qi + ˙Si ˙qi (a0 = −ag) f B i = Ii ai + vi ×∗Ii vi fi = f B i −f x i + X j∈µ(i) fj τi = ST i fi Modified Equations: vi = vλ(i) + Si ˙qi (v0 = 0) avp i = avp λ(i) + ˙Si ˙qi (avp 0 = −ag) f B i = Ii avp i + vi ×∗Ii vi fi = f B i −f x i − NJ X k=NB+1 eikf a k + X j∈µ(i) fj where eik =

+1 if i = s(k) −1 if i = p(k) 0 otherwise C′ i = ST i fi Table 8.2: Equations for calculating C −τ a Table 8.2 shows the basic equations of the recursive Newton-Euler algorithm, as they appear in Table 5.1, alongside the equations for calculating C −τ a. On comparing the originals with the modified equations, the following differences can be seen. 1. The term Si ¨qi has been dropped, as the acceleration variables are set to zero for calculating C. 2. The variable names ai and τi have been changed to avp i and C′ i (where C′ = C −τ a). This is a purely cosmetic change.1 3. A new term has been added to the joint force calculation, which subtracts the active loop-joint forces at the same point where the general external forces are subtracted. From the perspective of the spanning tree, the active loop-joint forces are indeed external forces, since they come from something that is not a part of the tree, and that is exactly how they are treated. The coefficients eik ensure that f a k gets subtracted only from fs(k), and that −f a k gets subtracted from fp(k). Table 8.3 shows the pseudocode for an algorithm to calculate C −τ a according to the equations in Table 8.2. The calculations are performed in body coordinates. The first and third loops in this algorithm closely resemble the corresponding code in Table 5.1, but the second loop is new. The last line in loop 1 initializes each vector fi to the expression f B i −f x i . The purpose of loop 2 is then to subtract the term PNJ k=NB+1 eikf a k from each fi before loop 3 starts to use them. The second line in loop 2 uses jcalc to calculate T a k , which is placed in the local variable T a. There is an assumption here that T a k is a constant, so that 1Strictly speaking avp i in these equations has the value ˙Ji ˙q−ag, rather than ˙Ji ˙q as defined in Eq. 8.13, but the constant offset ag cancels out in the subtraction in Eq. 8.15.

v0 = 0 a0 = −ag for i = 1 to NB do [XJ, Si, vJ, cJ] = jcalc(jtype(i), qi, ˙qi) iXλ(i) = XJ XT(i) if λ(i)̸ = 0 then iX0 = iXλ(i) λ(i)X0 end vi = iXλ(i) vλ(i) + vJ avp i = iXλ(i) avp λ(i) + cJ + vi × vJ fi = Ii avp i + vi ×∗Ii vi −iX∗ 0 f x i end ... ... for l = 1 to NL do k = l + NB [T a] = jcalc(jtype(k)) f a = XS(l)T T a τk fs(k) = fs(k) −f a fp(k) = fp(k) + p(k)X∗ 0 0X∗ s(k) f a end for i = NB to 1 do C′ i = ST i fi if λ(i)̸ = 0 then fi = fi + λ(i)X∗ i fi end end Table 8.3: Algorithm to calculate C −τ a jcalc only needs to know the type specifier for joint k in order to work out the correct value. If we did not make this assumption, then it would be necessary to work out the transform from the predecessor frame to the successor frame of joint k, so that it can be supplied as an argument to jcalc. The expression T aτk gives f a k in the successor-frame coordinates of joint k, which are the coordinates in which T a is supplied. The third line in loop 2 calculates f a k in successor-frame coordinates, transforms it to body s(k) coordinates, and stores the result in the local variable f a. The next two lines then subtract this quantity from fs(k) and add it to fp(k). Observe that f a must be transformed from s(k) to p(k) coordinates before the addition is carried out. 8.7 Algorithm for K and k Table 8.4 shows the pseudocode for an algorithm to calculate K via Eqs. 8.19 and 8.20, k via Eq. 8.15, and kstab via Eqs. 8.22 and 8.23. These calculations involve combining vectors and matrices that were originally obtained in various different coordinate systems. It is therefore necessary to apply coordinate transforms to at least some of these quantities before they can be combined. The strategy used in this algorithm is to transform everything to base coordinates and combine them there. This is the simplest strategy, but not necessarily the best. The local variables Xp and Xs contain the Pl¨ucker transforms from base coordinates to the predecessor and successor frames, respectively, of loop joint k; so Xp = p(k),kX0 and Xs = s(k),kX0. The expression Xs X−1 p is therefore the transform from the joint’s predecessor frame to its successor frame, which is the argument to jcalc in Eq. 8.23. The call to jcalc requests both the loop position

K = 0 for l = 1 to NL do k = l + NB Xp = XP(l) p(k)X0 Xs = XS(l) s(k)X0 vp = 0Xp(k) vp(k) vs = 0Xs(k) vs(k) ap = 0Xp(k) avp p(k) as = 0Xs(k) avp s(k) [δ, T ] = jcalc(jtype(k), Xs X−1 p ) T = XT s T kstab = 2α T T(vs −vp) + β2 δ ... ... kl = −T T(as −ap + vs × vp) −kstab i = p(k) j = s(k) while i̸ = j do if i > j then Kli = −T T 0Xi Si i = λ(i) else Klj = T T 0Xj Sj j = λ(j) end end end Table 8.4: Algorithm to calculate K, k and kstab error, δ, and the constraint-force subspace matrix, T , for loop joint k. The line immediately after this call transforms T from successor-frame coordinates to base coordinates; but no transform is applied to δ because it represents a quantity of the form T Td, where d ∈M6 is an infinitesimal displacement. This algorithm assumes that ˚ T = 0 and σ = ˙σ = 0 for every loop joint. Without this assumption, jcalc would need more arguments, and would be asked to return more quantities. The assumption ˚ T = 0 implies ˙T = vs ×∗T , which allows us to simplify the right-hand side of Eq. 8.15 as follows: kl = −T T k (avp s(k) −avp p(k)) −˙T T k (vs(k) −vp(k)) = −T T k (avp s(k) −avp p(k)) −(vs(k) ×∗Tk)T(vs(k) −vp(k)) = −T T k (avp s(k) −avp p(k)) + T T k vs(k) × (vs(k) −vp(k)) = −T T k (avp s(k) −avp p(k) + vs(k) × vp(k)) . (8.47) This is the expression used in the algorithm. K is calculated by first initializing the whole matrix to zero, and then calculating the nonzero submatrices. For each loop l, the submatrix Kli is nonzero for every tree joint i that lies on the path of the loop. These quantities are calculated in a while loop that employs two loop variables, i and j. These variables are initially set to p(k) and s(k), respectively, and then stepped back toward the base in an order that guarantees they meet at the nearest common ancestor of p(k) and s(k). At this point, every joint on the path of loop l is accounted for, and we have i = j, which is the termination condition for the loop. The algorithm presented in Table 8.4 potentially involves some duplication of calculations. For example, if joint i participates in more than one loop, then

Si will be transformed to base coordinates more than once. Likewise, if body i is the successor or predecessor of more than one loop joint, then vi and avp i will be transformed to base coordinates more than once, and the relevant transform (Xp or Xs) will be calculated more than once. This kind of duplication has a relatively small impact on the overall efficiency. Nevertheless, if maximum speed is required, then the code can be rewritten to avoid such duplications. The Rounding-Error Problem The algorithm’s use of base coordinates achieves two things: it keeps the algorithm simple, and it minimizes the number of transformation matrices that are needed. However, there is a price to be paid: most floating-point calculations involve rounding errors, and some of these errors are sensitive to the distance between the frame in which the calculations are performed and the locations of the objects to which those calculations apply. Therefore, if the base coordinate frame is far away from the kinematic loop, then there will be a loss of precision in the calculations. This phenomenon is not a problem for most rigid-body systems; but it can be a problem for a floating-base system if it travels a great distance from its fixed-base origin. (Think of an aircraft, or a spacecraft.) In this case, a simple remedy is to perform the calculations for K, k and kstab in floating-base coordinates instead of fixed-base coordinates. 8.8 Algorithm for G and g This section shows how to solve a linear equation of the form K x = k (8.48) to yield a solution of the form x = G y + g . (8.49) The method is applicable to the task of solving Eq. 8.35 to get Eq. 8.42. In Equations 8.48 and 8.49, K is an m × n matrix of rank r, k ∈range(K) (so Eq. 8.48 is consistent), G is an n × (n −r) full-rank matrix, and Eq. 8.49 produces every possible solution to Eq. 8.48. In fact, Eq. 8.49 defines a 1 : 1 mapping between y and the set of solutions to Eq. 8.48. If r = n then Eq. 8.49 simplifies to x = g, and g is unique. If r < n then G and g are not unique, and we seek any one pair of values that solves the problem. Given that rank(K) = r, it follows that K must contain an r × r invertible submatrix. Therefore, there exists a permutation of K in which the elements of this submatrix occupy the top left corner. We may therefore write K = Q1 K11 K12 K21 K22  Q2 ,

in which K11 is an r × r invertible matrix, and Q1 and Q2 are permutation matrices. Extending this permutation to the whole of Eq. 8.48 gives Q1  K11 K12 K21 K22   x1 x2  = Q1  k1 k2  , (8.50) where x1 x2  = Q2 x and Q1 k1 k2  = k . Again, given that rank(K) = r and that Eq. 8.48 is consistent, it follows that the bottom row in Eq. 8.50 is a linear multiple of the top row; so there must exist a matrix Y such that K21 = Y K11, K22 = Y K12 and k2 = Y k1. We can therefore rewrite Eq. 8.50 as Q1  K11 K12 Y K11 Y K12  x1 x2  = Q1  k1 Y k1  . By inspection, we can see that one possible solution to this equation is x1 x2  =  −K−1 11 K12 1  x2 +  K−1 11 k1 0  . We may therefore conclude that one solution to the original problem is G = QT 2  −K−1 11 K12 1  and g = QT 2  K−1 11 k1 0  . (8.51) This particular solution equates y with x2, which means that y is a subset of the elements of x. The expressions in Eq. 8.51 can be calculated using a standard Gaussian elimination algorithm, or an LU factorization, with the following modifications. 1. Complete pivoting must be used, with both row and column swaps. A record must be kept of the column swaps, as this data determines the value of Q2. 2. If r has not been supplied as an argument, then a rank test is needed to detect when the last of the independent rows has been processed. 3. The factorization process terminates once it has processed the last independent row. At this stage, the top-left r×r submatrix of K contains the upper-triangular matrix U, where LU = K11; the top-right r × (n −r) submatrix contains L−1K12; and the top r rows of k contain L−1k1. Back-substitution on the top r rows therefore yields K−1 11 K12 and K−1 11 k1. Various features can be added to this basic algorithm. For example, a mask can be supplied that identifies which columns of K correspond to joint variables that do not participate in any kinematic loop. These columns contain only zeros, and are therefore left unaltered by the factorization process. Another possibility is to allow the user to specify his own choice of independent variables. This eliminates the need for complete pivoting, so that partial pivoting can be used instead.

k1 k2 k3 a b c d e f K1 K3 K2 g1 g2 g3 a b c d e f G1 1 1 1 G2 G3 K k G g 1 2 3 4 5 6 7 8 9 10 11 17 16 15 14 13 12 20 19 18 a b c d e f 24 23 22 21 Figure 8.1: Loop clusters and the sparsity patterns of K and G 8.9 Exploiting Sparsity in K and G It is clearly the case that both K and G may contain many zeros. For example, Eq. 8.20 implies that each submatrix Klj of K will be zero for every l and j such that εlj = 0; and Eq. 8.51 implies that G contains at least (n−r)(n−r−1) zeros, this being the number of zeros in an (n −r) × (n −r) identity matrix. However, a large proportion of all the sparsity in these two matrices can be attributed to a single fact, which is the subject of this section: they are both block-diagonal. To reveal the block-diagonal structure, we must first organize the kinematic loops into a set of minimal loop clusters. A loop cluster is any set of loops with the property that no loop within the cluster is coupled to any loop outside it; and a minimal loop cluster is one that cannot be divided into two smaller clusters. A coupling exists between two kinematic loops if they have a joint in common. Given a set of minimal loop clusters, it is possible to devise a regular numbering of the spanning tree with the following property: the tree joints within each cluster are numbered consecutively. It is this numbering scheme that reveals the block-diagonal structure of K and G. An example is presented in Figure 8.1. This figure shows the connectivity graph of a closed-loop system containing four kinematic loops. They are numbered 1 to 4, and are closed by loop joints 21 to 24, respectively. (Loop numbers

are not shown in the diagram.) Loop 1 is coupled to loop 2, but there is no other coupling in the system. We can therefore identify three minimal clusters as follows: cluster 1 contains loops 1 and 2; cluster 2 contains loop 3 only; and cluster 3 contains loop 4 only. Having identified these three clusters, we can devise a regular numbering in which the tree joints are numbered consecutively within each cluster. In this example, the joints in cluster 1 are numbered 3 to 7, those in cluster 2 are numbered 10 to 12, and those in cluster 3 are numbered 19 and 20. The remaining tree joints do not lie on the path of any kinematic loop, and therefore do not belong to any cluster. They can be numbered in any manner that is consistent with the rules of regular numbering. With this numbering scheme, K adopts an irregular block-diagonal structure with gaps: there is one block per cluster; blocks are generally rectangular; each block can be a different size and shape to the others; and there can be horizontal gaps of varying sizes between the blocks. These gaps are caused by the tree joints that do not belong to any cluster. This pattern of sparsity in K allows G to adopt an irregular block-diagonal structure without gaps: each block Ki in K gives rise to a block Gi in G; and each gap in K gives rise to an identity-matrix block in G. For easy identification, the joints in the connectivity graph have been divided into six sets, labelled a to f, and the same division is shown against the columns of K and the rows of G. Block-diagonal sparsity can be exploited in the calculation of KH−1KT and GTHG, and in the calculation of G from K. In particular, G can be calculated from K one block at a time: as x = Gi y + gi is the solution to Ki x = ki, each Gi and gi can be calculated directly from Ki and ki using the algorithm in Section 8.8. 8.10 Some Properties of Closed-Loop Systems Closed-loop systems exhibit a greater variety of properties and behaviours than kinematic trees, which can pose new difficulties for dynamics algorithms. This section looks at some of the special properties of closed-loop systems that are relevant to dynamics. Mobility: Mobility is the degree of motion freedom of a rigid-body system. The mobility of a kinematic tree is simply n; but the mobility of a closed-loop system is given by the general formula mobility = n −r . (8.52) If y is a vector of independent coordinates for a rigid-body system, then the dimension of y is n −r. Systems with the property r = nc are said to be properly constrained. For such a system, the mobility can be expressed by the formula mobility = ntot −6NL , (8.53)

θ (a) θ θ ? (b) Figure 8.2: Examples of variable mobility (a) and configuration ambiguity (b) where ntot is the sum of the joint freedoms of all the joints in the system, not just the tree joints. This formula applies to rigid-body systems in a 3D Euclidean space. The equivalent formula in a 2D space is mobility = ntot −3NL . (8.54) Variable Mobility: The mobility of a kinematic tree is fixed; but the mobility of a closed-loop system depends on r, and r can vary with configuration. An example of variable mobility is shown in Figure 8.2(a). As long as the angle θ is nonzero, the mechanism has a mobility of 1; but if θ = 0 then the two arms are able to move independently, and the mobility is 2. Furthermore, at the configuration shown in the middle diagram, which is the boundary between these two motion regimes, the mechanism has an instantaneous mobility of 3. In terms of n and r, we have n = 5 throughout, but r = 4 in the left-hand diagram, r = 3 in the right-hand diagram, and r = 2 in the middle diagram. Configuration Ambiguity: If q and y are vectors of independent position variables for a kinematic tree and a closed-loop system, respectively, then q always identifies a unique configuration of the tree, but y does not necessarily identify a unique configuration of the closed-loop system. A simple example is shown in Figure 8.2(b): knowing the value of the independent variable, θ, is not enough to determine uniquely the configuration of the system. There are two ways to solve this problem: 1. use q instead of y, where q is the position vector of the spanning tree, or 2. supply a function that maps each value of y to a unique value of q.

The first option is the one used in previous sections, while the second is the subject of Section 8.11. For the system shown in Figure 8.2(b), option 2 is applicable if we know in advance that one of the two configurations cannot be reached. Otherwise, we must either try a different independent variable, or else resort to option 1. Overconstraint: A closed-loop system is said to be overconstrained if r < nc. Such systems have nc −r redundant constraints, which means there will be nc −r indeterminate constraint forces in the equations of motion. They also have excess mobility relative to a properly constrained system, as the mobility formula can be written mobility = ntot −6NL + (nc −r) . (8.55) Overconstrained systems are actually very common. For example, any system containing planar kinematic loops will be overconstrained. The planar loop shown in Figure 8.2(b) contains four revolute joints, so the numbers for this system are n = 3 and nc = 5; but the mobility is 1, which means r = 2. The dynamics of overconstrained systems is harder to calculate than that of properly constrained systems. For this reason, it is useful to convert the former to the latter wherever possible. The conversion is achieved by replacing the original loop joints with joints that impose fewer constraints, so that the value of nc is reduced. For example, if the loop-closing revolute joint in a planar kinematic loop is replaced with a sphere-in-cylinder joint, which imposes only two constraints, then nc will be reduced from 5 to 2, which means nc = r. 8.11 Loop Closure Functions Let q be the vector of position variables for the spanning tree of a given closedloop system, and let y be a vector of independent position variables for the same system. y will typically be a subset of the elements of q, but this is not necessary. We assume that y defines q uniquely. For this assumption to hold, it may be necessary to place restrictions on the value of y. We therefore define a set C ⊆R(n−r) of acceptable values for y, and assume y ∈C. Given these assumptions, there exists a function, γ, that satisfies q = γ(y) (8.56) for all y ∈C. (Cf. Eq. 3.11.) Differentiating this equation gives ˙q = G ˙y , (8.57) where G = ∂γ ∂y , (8.58)

and differentiating it a second time gives ¨q = G ¨y + g , (8.59) where g = ˙G ˙y . (8.60) Equation 8.59 has the same form as Eq. 8.42, and the two equations are identical when the loop-closure errors are zero. We now have a new method for obtaining G and g, in which they are calculated from a vector of independent variables. This method can be summarized as follows. 1. Define y, and identify C. 2. Find an expression for the function γ that maps y uniquely to q, given y ∈C. 3. Differentiate this expression symbolically to obtain expressions for G and g. 4. Write code to evaluate these expressions. This method is less general that the one described in Section 8.5, mainly because step 2 is not always possible. It is also less convenient, since it requires the user to supply expressions (or code) for G and g in addition to the system model. Nevertheless, it is applicable to a wide variety of practical closed-loop systems, and the calculation of G and g by this method can be much cheaper than using the algorithms in Sections 8.7 and 8.8. Another advantage of this method is that loop-closure errors cannot occur. This is because we do not integrate ¨q to get ˙q and q, but instead integrate ¨y to get ˙y and y. This way, the vectors q and ˙q are calculated from y and ˙y using Eqs. 8.56 and 8.57, which means that their values automatically satisfy the loop constraints. There is therefore no need for constraint stabilization terms. A Worked Example Figure 8.3 shows a rigid-body system containing four bodies, five joints and one kinematic loop. The system is a planar mechanism, and its mobility is 2. Bodies B1 and B4 form a two-link arm, while B2 and B3 are the cylinder and piston, respectively, of a linear actuator. Joint 3 is prismatic, and the rest are revolute. Joints 1, 2, 3 and 5 participate in the kinematic loop, with joint 5 being the loop joint. The vector of tree-joint variables is therefore q = [q1 q2 q3 q4]T, where q3 is a distance and the rest are angles. In this example, we choose q1 and q4 to be the independent variables; that is, we define the vector of independent variables, y = [y1 y2]T, such that y1 = q1 and y2 = q4. The main task is then to express q2 and q3 as functions of y1.

θ1 (=y1) θ2 (=q2) l1 l2 d (=q3) B1 B4 q1 (=y1) q2 q4 (=y2) q3 l1 l2 B3 B2 loop geometry J3 Triangle Formulae: d2 = l2 1 + l2 2 −2l1l2 cos(θ1) d > 0 sin(θ2) = l2 sin(θ1) d cos(θ2) = l2 cos(θ1) −l1 d Solution: γ(y) =

y1 θ2 d y2

where θ2 = tan−1  l2 sin(y1) l2 cos(y1) −l1  d = q l2 1 + l2 2 −2l1l2 cos(y1) G = ∂γ ∂y =

1 0 u1 0 u2 0 0 1

where u1 = l2 2 −l1l2 cos(y1) d2 u2 = l1l2 sin(y1) d g = ˙G ˙y =

0 u3 u4 0

where u3 = (l2 1 −l2 2)l1l2 sin(y1) d4 ˙y2 1 u4 = (l1l2 cos(y1) −u2 2) d ˙y2 1 Figure 8.3: Formulating γ, G and g for a simple closed-loop system

The geometry of the kinematic loop is very simple: just a triangle. It is therefore possible to use standard trigonometric formulae to express q2 and q3 (labelled θ2 and d on the triangle) as functions of y1. To make the mapping unique, we impose the constraint d > 0, and we use the four-quadrant version of the arctangent function (the one that takes two arguments) to calculate θ2. If l1 = l2 then we require θ1̸ = 0. In this case, C must exclude all vectors y having y1 = 0. When a kinematic loop has an explicit solution formula, like this one, then it is said to have a closed-form solution. Figure 8.3 also presents explicit formulae for γ, G and g. It can be seen that G is block-diagonal, having a 3 × 1 block and a 1 × 1 block. Note that the equations in this figure are all scalar equations, in contrast to the vector and matrix equations in Sections 8.7 and 8.8. It is therefore much cheaper to calculate G and g via the equations presented here. In summary, this method places an additional burden on the user of having to supply code to calculate γ, G and g, but offers in return a potentially large improvement in the efficiency of this part of the dynamics calculation. This method is appropriate for systems containing kinematic loops with simple, closed-form solutions, such as might be found on excavator arms, large industrial robots, parallel robots and Stuart-Gough platforms (flight simulators). 8.12 Inverse Dynamics Inverse dynamics is the problem of calculating the forces required to produce a given acceleration. In a kinematic tree, there is a 1 : 1 relationship between τ and ¨q, and the calculation of τ from ¨q is straightforward. In a closed-loop system, the given value of ¨q must comply with the loop-closure constraints, and there are infinitely many values of τ that produce the same acceleration. The nature of the inverse dynamics problem is therefore significantly altered by the presence of kinematic loops. Let us now formulate the inverse dynamics equations for a closed-loop system. We start with the equation of motion for the spanning tree, which is τ = H¨q + C −τ a −KTλ . (8.61) In this equation, τ is the vector of unknown forces acting at the tree joints. Any known forces, such as those due to springs and dampers, are assumed to be included in C. It is also assumed that there are no unknown active forces at the loop joints. Thus, if τ a appears at all, then it is composed entirely of known forces. Equation 8.61 contains two unknowns, τ and λ. As we are only interested in τ, the next step is to eliminate λ. This is done by multiplying the equation by GT, which gives GTτ = GT(H¨q + C −τ a) .

This equation can be written GTτ = GTτID , (8.62) where τID is the force calculated by applying an inverse dynamics function to the spanning tree; that is, τID = H¨q + C −τ a = ID(q, ˙q, ¨q) . In this context, ID(q, ˙q, ¨q) is the output of the algorithm described in Section 8.6 with the joint acceleration terms put back in. Bear in mind that q, ˙q and ¨q must all comply with the loop-closure constraints. Equation 8.62 is the basic inverse dynamics equation for a closed-loop system. As GT is an (n −r) × n matrix, this equation imposes n −r constraints on an n-dimensional vector of unknowns, leaving r freedoms of choice. In other words, there are ∞r different values of τ that produce the same acceleration. If we want a unique solution to the problem, then we must either impose more constraints or apply an optimality criterion. Actuator Forces A typical use for inverse dynamics is to calculate a vector of actuator forces to produce a given acceleration. Now, not every joint will be powered by an actuator, so let us define p and u to be the number of actuated freedoms in the tree, and the p-dimensional vector of actuator force values, respectively. u and τ are then related by the equation τ = Q u 0  , where Q is an n × n permutation matrix. This equation simply says that the elements of τ are a permutation of the p elements of u and n −p zeros. The zero-valued elements of τ are those that belong to the unactuated degrees of freedom. Let us also define the two matrices Gu and G0, which contain the rows of G corresponding to actuated and unactuated freedoms, respectively. These matrices are related to G by the equation G = Q Gu G0  . From these definitions, it follows that GTτ =  GT u GT 0  QT Q  u 0  =  GT u GT 0   u 0  = GT uu ; and substituting this result into Eq. 8.62 gives GT uu = GTτID , (8.63)

which is the inverse dynamics equation for actuator forces. In solving this equation, there are three cases to consider. 1. rank(Gu) < n −r. In this case, the system is underactuated: it has more motion freedoms than the actuators can control. 2. p > rank(Gu). In this case, the system is redundantly actuated: infinitely many different values of u will produce the same acceleration. 3. p = rank(Gu) = n −r. In this case, the system is properly actuated: Gu is invertible, and there is a unique solution to Eq. 8.63 given by u = G−T u GTτID. Cases 1 and 2 are not mutually exclusive. Example 8.2 Consider the assignment of actuators to the mechanism shown in Figure 8.3. This mechanism has two degrees of freedom, and therefore requires two actuators. Gu will therefore be a 2 × 2 matrix containing two rows of G; and these rows must be chosen so that Gu is invertible. Given the formula for G shown in the figure, it follows that Gu must contain row 4 of G and any one of the other three rows. This means that an actuator must be present at joint 4, and at any one of the other three joints. (As this mechanism is so simple, this result could have been obtained by inspection.) Let us suppose that joints 3 and 4 are actuated, which means that Gu =  u2 0 0 1  . The inverse dynamics formula for this system will then be u = G−T u GT τID =  u2 0 0 1 −1  1 u1 u2 0 0 0 0 1  τID = 1/u2 u1/u2 1 0 0 0 0 1  τID , where τID = ID(γ(y), G ˙y, G¨y + g) . Note that u is not a vector of generalized forces for this system, since the expression uT ˙y does not equal the power delivered by u. This discrepancy is caused by the fact that ˙q1 is the first element in ˙y, but τ3 is the first element in u. If you want u to be the generalized force vector, then the elements of u must be the same subset of τ as ˙y is of ˙q. 8.13 Sparse Matrix Method Every algorithm we have studied so far works by applying a set of loop-closure constraints to a spanning tree. An alternative strategy is to dispense with the spanning tree, and work directly with the individual bodies and joints in the

system. This approach produces equations with large, sparse matrices; and any algorithm based on these equations must exploit the sparsity in order to be competitive. Such algorithms tend to be less efficient than those based on spanning trees, but they can be more efficient under special circumstances, as we shall see below. The equation of motion for a collection of NB rigid bodies, subject to the motion constraints imposed by a set of NJ joints, is given by  I P T T TP T 0   a −λ  = P T aτ −v ×∗Iv + f x −˙T TP Tv  . (8.64) This equation is derived in Section 3.7. v, a and f x are vectors of body velocities, accelerations and external forces, respectively; λ is the vector of joint constraint-force variables; T aτ is the vector of active forces transmitted across the joints; P T is the matrix that maps from body to joint velocities (i.e., P Tv is the vector of joint velocities); and I and T are composite matrices of body inertias and joint constraint-force subspaces, respectively. a and λ are the unknowns; and the size of the coefficient matrix is 6NB + nc tot square, where nc tot is the total number of constraints imposed by the joints. Our chief interest lies in the sparsity pattern of the coefficient matrix, so let us examine its constituents. I and T are both block-diagonal, being defined by I = diag(I1, . . . , INB) and T = diag(T1, . . . , TNJ) . P is also sparse, but not block-diagonal. It is defined by P ij =

16×6 if i = s(j) −16×6 if i = p(j) 06×6 otherwise. P is therefore a 6NB × 6NJ matrix, organized as an NB × NJ matrix of 6 × 6 blocks. It has two nonzero blocks in every column, except those columns for joints connecting directly to the base, which have only one nonzero block. From these definitions, it follows that the coefficient matrix contains NB nonzero blocks on the main diagonal, and up to 4NJ nonzero off-diagonal blocks. The exact positions of the latter are determined by the connectivity. Equation 8.64 is typical of the kind of equation that a sparse matrix arithmetic software library is designed to solve, and the use of such a library is probably the best way to proceed. A sparse matrix library would typically provide the following facilities: • a function to analyse the sparsity pattern of a matrix, as a preprocessing step, in order to work out the best factorization order and which zerovalued elements get filled in with nonzero values during the course of the factorization; • a data structure for compact storage of all the nonzero values in the matrix and all the zero-valued entries that will get altered during factorization;

• a function to perform a factorization in the order determined in the preprocessing step; and • a function to perform back-substitution using the resulting factors. If an n × n matrix contains only O(n) nonzero elements, then it is sometimes possible to factorize the matrix in only O(n) arithmetic operations. An O(n) factorization of Eq. 8.64 is always possible if the system is a kinematic tree.2 However, there are also many closed-loop systems for which an O(n) factorization is possible, including some that cannot be solved in O(n) by any algorithm using a spanning tree. In such cases, the sparse-matrix algorithm will out-perform spanning-tree algorithms for a sufficiently large n. Figure 8.4 shows an example of a closed-loop system having O(n) dynamics by the sparse matrix method. The connectivity graph is of a kind that is sometimes called a ladder. It can be extended indefinitely by adding more nodes to the right. The figure also shows the sparsity pattern of a symmetric permutation of the coefficient matrix. This permutation has been chosen so that the resulting matrix has a fixed bandwidth—every nonzero block appears on the main diagonal, or on one of the three nearest diagonals either side. A matrix with this property can always be factorized in O(n) operations. To explain the permutation, the nodes in the graph have been numbered in the same order as they appear in the matrix, and the joints have been given number pairs reflecting their connectivity: joint ij is the one that connects from body i to body j. Row 1 in the permuted matrix is the row in the original matrix that contained I1 (with body numbers as shown in the graph); row 2 is the row that contained I2; row 3, which is labelled 12, is the row in the original matrix containing the row in T TP T pertaining to joint 12; and so on. Observe that each row labelled with a body number contains a nonzero block on the main diagonal, and that each row labelled with a number pair ij contains nonzero blocks only in columns i and j. If we were to apply a spanning-tree algorithm to this system, then we would encounter the following problem: whatever spanning tree we choose, there will always be O(n2) nonzero elements in the joint-space inertia matrix, H, and so it will be impossible to calculate the dynamics in only O(n) operations. Finally, the following details should be mentioned. First, if this algorithm is to be used in a dynamics simulator, then constraint stabilization is required. The necessary extra terms have not been included in Eq. 8.64. Second, we have assumed that the coefficient matrix has full rank; that is, we have assumed that the closed-loop system is properly constrained. If the matrix is singular, then the sparse factorization may not work. 2This is the basis of Baraff’s algorithm (Baraff, 1996).

1 2 12 13 24 3 4 34 35 46 5 6 56 57 68 7 8 78 70 80 1 2 12 13 24 3 4 34 35 46 5 6 56 57 68 7 8 78 70 80 1 2 3 4 5 6 7 8 0 12 34 56 78 13 24 46 35 57 68 70 80 Figure 8.4: Sparsity pattern arising from a ladder-like connectivity graph
